% ============================================================
% preamble.tex — configuración compartida, compatible con IEEEtran
%
% NO agregar: geometry, fancyhdr, setspace, tocloft, natbib,
% caption, subcaption, cleveref. IEEEtran los maneja internamente.
% ============================================================

% ---------- Idioma ----------
\usepackage[T1]{fontenc}
\usepackage[spanish,es-tabla,es-noquoting,es-nodecimaldot]{babel}
\usepackage{microtype}

% ---------- Encabezados de IEEEtran en español ----------
\renewcommand{\abstractname}{Resumen}
\renewcommand{\IEEEkeywordsname}{Palabras clave}
\renewcommand{\refname}{Referencias}
\renewcommand{\figurename}{Fig.}
\renewcommand{\tablename}{TABLA}

% ---------- Numeración de página ----------
% IEEEtran conference la desactiva por defecto porque en actas
% reales la asigna el editor. Se activa por ser trabajo de curso.
\pagestyle{plain}
\AtBeginDocument{\thispagestyle{plain}}

\usepackage{amsmath}
\usepackage{amssymb}
\usepackage{amsthm}
\usepackage{mathtools}
\usepackage{siunitx}
\sisetup{per-mode=symbol}

\usepackage{cite}

\usepackage{xcolor}
\definecolor{tecnavy}{RGB}{31,56,100}
\definecolor{tecsteel}{RGB}{46,85,150}
\definecolor{tecrule}{RGB}{155,184,216}
\definecolor{boxbg}{RGB}{242,246,252}
\definecolor{labelgray}{RGB}{120,120,120}
\definecolor{notaverde}{RGB}{22,110,60}
\definecolor{decisionambar}{RGB}{170,110,10}
\definecolor{codebg}{RGB}{248,248,248}
\definecolor{codekw}{RGB}{0,0,180}
\definecolor{codecmt}{RGB}{0,120,60}
\definecolor{codestr}{RGB}{170,55,55}

\usepackage{graphicx}
\graphicspath{{img/}{figures/}{./}}
\usepackage{tikz}
\usepackage{pgfplots}
\pgfplotsset{compat=1.18}
\usetikzlibrary{arrows.meta,positioning,shapes.geometric,calc,fit,
                backgrounds,decorations.pathreplacing}

\usepackage{array}
\usepackage{booktabs}
\usepackage{multirow}
\usepackage{makecell}
\usepackage{tabularx}
\usepackage{threeparttable}
\usepackage{enumitem}

\newcolumntype{Y}{>{\raggedright\arraybackslash}X}
\newcolumntype{P}[1]{>{\raggedright\arraybackslash}p{#1}}
\renewcommand{\arraystretch}{1.15}
\setlist[itemize]{leftmargin=*,noitemsep,topsep=2pt}
\setlist[enumerate]{leftmargin=*,noitemsep,topsep=2pt}

\theoremstyle{definition}
\newtheorem{definicion}{Definición}
\newtheorem{ejemplo}{Ejemplo}
\theoremstyle{plain}
\newtheorem{proposicion}{Proposición}
\theoremstyle{remark}
\newtheorem*{observacion}{Observación}

\usepackage{listings}
\lstdefinestyle{code}{
  basicstyle=\ttfamily\scriptsize,
  keywordstyle=\color{codekw}\bfseries,
  commentstyle=\color{codecmt}\itshape,
  stringstyle=\color{codestr},
  backgroundcolor=\color{codebg},
  numbers=left,
  numberstyle=\tiny\color{labelgray},
  numbersep=6pt,
  frame=single,
  rulecolor=\color{tecrule},
  breaklines=true,
  postbreak=\mbox{\textcolor{tecsteel}{$\hookrightarrow$}\space},
  showstringspaces=false,
  tabsize=2,
  captionpos=b,
  columns=flexible,
  upquote=true
}
\lstset{style=code}

\usepackage[hidelinks]{hyperref}

% ---------- Entrada segura ----------
\newcommand{\cajaarchivofaltante}[1]{%
  \par\medskip\noindent
  \fcolorbox{red!60!black}{red!5}{%
    \parbox{\dimexpr\linewidth-2\fboxsep-2\fboxrule}{%
      \centering\ttfamily\scriptsize ARCHIVO FALTANTE: \detokenize{#1}}}%
  \par\medskip}

\newcommand{\cajaimagenfaltante}[1]{%
  \begingroup\setlength{\fboxsep}{0pt}%
  \fcolorbox{tecrule}{boxbg}{%
    \parbox[c][2.4cm][c]{0.9\linewidth}{\centering\ttfamily\scriptsize
      IMAGEN NO ENCONTRADA\\[3pt]\detokenize{#1}}}%
  \endgroup}

\newcommand{\safeinput}[1]{%
  \IfFileExists{#1.tex}{\input{#1}}{%
    \PackageWarning{tecpaper}{Falta el archivo de sección #1.tex}%
    \cajaarchivofaltante{#1.tex}}}

\newcommand{\safeimage}[2][width=\columnwidth]{%
  \IfFileExists{#2}{\includegraphics[#1]{#2}}{%
    \IfFileExists{img/#2}{\includegraphics[#1]{img/#2}}{%
      \PackageWarning{tecpaper}{Falta la imagen #2}%
      \cajaimagenfaltante{#2}}}}

\newcommand{\code}[1]{\texttt{#1}}

% ============================================================
% MARCADORES DE REDACCIÓN
%
% Los tres se apagan antes de entregar redefiniéndolos como vacíos:
%   \renewcommand{\pc}[1]{}
%   \renewcommand{\razon}[1]{}
%   \renewcommand{\decidir}[1]{}
% Recompilar: todo lo que desaparezca del PDF nunca se escribió.
% ============================================================

% \pc  = POR COMPLETAR. Contenido que depende de los ejes A y C.
\newcommand{\pc}[1]{%
  \texorpdfstring{\textcolor{red}{\textbf{[POR COMPLETAR: #1]}}}%
                 {[POR COMPLETAR: #1]}}

% \razon = POR QUÉ esta parte debe decir lo que dice, con su fuente.
\newcommand{\razon}[1]{%
  \texorpdfstring{\textcolor{notaverde}{\small\textit{[RAZÓN: #1]}}}%
                 {[RAZÓN: #1]}}

% \decidir = requiere decisión de equipo o criterio humano.
\newcommand{\decidir}[1]{%
  \texorpdfstring{\textcolor{decisionambar}{\small\textbf{[DECIDIR: #1]}}}%
                 {[DECIDIR: #1]}}

% ---------- Columnas X con peso, para tabularx ----------
% Permite repartir el ancho flexible en proporciones distintas.
% Los pesos de todas las columnas W deben sumar el número de
% columnas W de la tabla.
\newcolumntype{W}[1]{>{\hsize=#1\hsize\raggedright\arraybackslash}X}

% ---------- Marcador compacto para celdas de tabla ----------
% \pc dentro de una celda produce una cadena larga que rompe el
% ajuste de columnas. Este marcador ocupa tres caracteres.
% Se apaga igual que los demás: \renewcommand{\pct}{}
\newcommand{\pct}{\textcolor{red}{\textbf{[?]}}}

% ============================================================
% APAGADO DE MARCADORES
% Descomentar las cuatro líneas antes de entregar y recompilar.
% Todo lo que desaparezca del PDF nunca se escribió.
% ============================================================

\renewcommand{\pc}[1]{}
\renewcommand{\razon}[1]{}
\renewcommand{\decidir}[1]{}
\renewcommand{\pct}{}